\expandafter\def\csname HistoricalNoteText1\endcsname{Unter einer imaginären Größe verstehe ich hier stets eine Größe der Form $a+b\sqrt{-1}$, solange $b$ nicht $=0$ ist. In diesem Sinn haben alle Mathematiker ersten Ranges jenen Ausdruck stets gebraucht. Ich meine, man sollte diejenigen nicht anhören, die eine Größe $a+b\sqrt{-1}$ nur dann imaginär nennen wollen, wenn $a=0$ ist, dagegen unmöglich, wenn nicht $a=0$ ist; denn diese Unterscheidung ist weder nötig noch von irgendeinem Nutzen. --- Sollen imaginäre Größen überhaupt in der Analysis beibehalten werden (was aus mehreren Gründen ratsamer erscheint, als sie abzuschaffen, sofern sie nur hinreichend sicher begründet werden), so muss man sie notwendig als ebenso möglich wie die reellen betrachten. Daher möchte ich die reellen und die imaginären Größen lieber unter der gemeinsamen Bezeichnung \textit{mögliche Größen} zusammenfassen. Dagegen würde ich eine Größe \textit{unmöglich} nennen, die Bedingungen erfüllen soll, die sich selbst bei Zulassung imaginärer Größen nicht erfüllen lassen, jedoch so, dass diese \textit{Ausdrucksweise} dasselbe bedeutet, wie wenn man sagt, eine solche Größe gebe es im gesamten Bereich der Größen nicht. Daraus eine besondere Art von Größen zu bilden, würde ich keineswegs gestatten. Sagt jemand, ein geradliniges gleichseitiges rechtwinkliges Dreieck sei unmöglich, so wird niemand widersprechen. Wollte er aber ein solches unmögliches Dreieck als eine neue Art von Dreiecken betrachten und andere Eigenschaften der Dreiecke darauf anwenden, wer könnte das Lachen zurückhalten? Das hieße mit Worten spielen oder vielmehr sie missbrauchen. --- Zwar haben auch die größten Mathematiker häufig Wahrheiten, die die Möglichkeit der Größen, auf die sie sich beziehen, offenkundig voraussetzen, auf solche angewendet, deren Möglichkeit noch zweifelhaft war. Auch will ich nicht bestreiten, dass derartige Freiheiten meist nur die Form und gleichsam die Hülle der Schlussfolgerungen betreffen, die der Scharfblick eines wahren Mathematikers bald durchdringen kann. Dennoch scheint es ratsamer und der Erhabenheit einer Wissenschaft, die mit Recht als vollkommenstes Muster von Klarheit und Gewissheit gerühmt wird, angemessener, solche Freiheiten entweder ganz zu verbannen oder sie wenigstens sparsamer und nur dort zu gebrauchen, wo auch weniger Geübte erkennen können, dass die Sache ohne ihre Hilfe vielleicht weniger kurz, aber ebenso streng hätte erledigt werden können. --- Im Übrigen will ich nicht bestreiten, dass man das, was ich hier gegen den Missbrauch der unmöglichen Größen gesagt habe, in gewisser Hinsicht auch gegen die imaginären einwenden könnte; doch behalte ich mir deren Verteidigung sowie eine ausführlichere Darstellung dieser ganzen Sache für eine andere Gelegenheit vor.}

